\beginsong{Le vent des prophètes}[by={Pierre-Michel Gambarelli}]
\beginverse
Tu \[Am]as fait se lever un peuple \[E]sans levain,
Le f\[Am]roment est battu à la for\[E]ce des poings.
Le \[G]pain cuit au soleil a pris le \[C]goût du sang,
Ton \[Dm]nom vient prolonger la liste \[E7]des gênants.
\endverse
\beginchorus
\[Am]Mais le \[Am]vent des pro\[G]phètes a souf\[F]flé \[G]ce ma\[E7]tin,
Et l'on a \[Am]vu des milliers d'alou\[G]ettes
Dan\[F]ser aut\[Em]our du pèle\[Am]rin.
Et l'on a \[Am]vu sur toute la pla\[G]nète,
Des \[F]frères \[Em]se donner la \[Am]main.
\endchorus
\beginverse
L'\[Am]injustice ne peut supporter \[E]le silence,
Mais \[Am]l'exemple du Christ prend toute \[E]sa puissance.
Sus\[G]pendu au gibet le jour de l'\[C]abandon,
Pour \[Dm]notre humanité il demand\[E7]e pardon.
\endverse
\beginverse
Tu \[Am]es mort le printemps n'avait que \[E]quelques brins,
Tué \[Am]dans ton église, les statues \[E]pour témoins.
Prêt\[G]re d'El Salvador, libre à \[C]en mourir,
Ton \[Dm]nom vient prolonger la liste \[E7]des martyrs.
\endverse
\beginverse
Com\[Am]bien d'hommes sont morts et combien \[E]meurent encore,
Ils \[Am]sont pour notre monde cet envers \[E]du décor.
Cach\[G]és par les parades et les marches \[C]militaires,
Qui \[Dm]pourrait désormais les forcer \[E7]à se taire ?
\endverse
\endsong
